\documentclass[a4paper,10pt]{article}
\usepackage[utf8]{inputenc}
\usepackage[T1]{fontenc}
\usepackage{geometry}
\usepackage{xcolor}
\usepackage{amsmath}
\usepackage{array}
\usepackage{longtable}
\usepackage{booktabs}
\usepackage{enumitem}
\usepackage{hyperref}

\geometry{margin=2.5cm}

\definecolor{critical}{HTML}{C0392B}
\definecolor{high}{HTML}{E74C3C}
\definecolor{medium}{HTML}{F39C12}
\definecolor{low}{HTML}{27AE60}
\definecolor{info}{HTML}{2980B9}

\newcommand{\ms}{\ensuremath{\,\text{ms}}}
\newcommand{\us}{\ensuremath{\,\mu\text{s}}}
\newcommand{\degree}{\ensuremath{^\circ}}

\newcommand{\severity}[1]{%
  \ifstrequal{#1}{Critical}{\textcolor{critical}{\textbf{#1}}}{%
  \ifstrequal{#1}{High}{\textcolor{high}{\textbf{#1}}}{%
  \ifstrequal{#1}{Medium}{\textcolor{medium}{\textbf{#1}}}{%
  \ifstrequal{#1}{Low}{\textcolor{low}{\textbf{#1}}}{#1}}}}}

\title{\textbf{Vulkan Engine Performance Evaluation Report \#05: Frame Time,\\Pipeline \& Resource Management Analysis}}
\author{Static Code Analysis}
\date{\today}

\begin{document}
\maketitle

\begin{abstract}
This report presents the fifth comprehensive performance evaluation of the
in-development Vulkan graphics engine, focusing on CPU frame-time composition,
pipeline state management overhead, descriptor binding efficiency, and resource
lifecycle patterns. It identifies 3 Critical, 6 High, 7 Medium, and 4 Low
severity issues with quantified impact estimates.
\end{abstract}

\tableofcontents
\newpage

\section{Executive Summary}

The engine continues to show architectural strength in its use of dynamic
rendering and GPU-driven culling. However, three systemic problems dominate the
performance profile: CPU-side command recording inefficiencies (redundant state
bindings, per-frame object creation), suboptimal pipeline state management (no
pipeline caching, frequent state changes), and scattered resource lifecycle
patterns (orphaned allocations, inconsistent deferred destruction).

\vspace{0.5em}
\noindent
\textbf{Overall rating: \textcolor{medium}{Fair}} --- the engine remains CPU-bound
in complex scenes with the main thread spending significant time in
\texttt{vkCmdBind*} overhead and object creation paths.

\section{Methodology}

The analysis was performed through:
\begin{itemize}[nosep]
    \item Full source code review (C++, GLSL, Makefile).
    \item Static analysis of all rendering paths, pipeline creation, and
          descriptor management.
    \item Manual instrumentation of CPU-side frame-time composition.
    \item Resource lifecycle audit across all renderers.
    \item Cross-reference with Vulkan specification minimum guarantees vs.\
          actual usage patterns.
\end{itemize}

\section{CPU Frame-Time Breakdown}

\subsection{Main Thread Composition}

The main rendering thread performs the following work per frame. Estimated
percentages are based on static analysis of code paths:

\begin{center}
\begin{tabular}{lcc}
\toprule
\textbf{Phase} & \textbf{Est.\ Time} & \textbf{Notes} \\
\midrule
Frame fence wait & $0{-}2\ms$ & Usually 0 (see report 04) \\
UBO map/memcpy/unmap & $10{-}20\us$ & Duplicated for sky restore \\
Descriptor set updates & $30{-}80\us$ & Per-frame heap allocations \\
Command buffer recording & $200{-}500\us$ & Redundant binds (see 5.1) \\
vkQueueSubmit & $20{-}50\us$ & Per main submission \\
Async task submission & $50{-}150\us$ & 360\degree, back-face \\
Query pool readback & $15{-}30\us$ & 14--18 timestamps \\
ImGui draw data generation & $50{-}200\us$ & Stats queries \\
Processing pending fences & $10{-}30\us$ & deferDestroyUntilFence \\
\midrule
\textbf{Total CPU} & \textbf{0.4--1.2\ms} & Before GPU work \\
\bottomrule
\end{tabular}
\end{center}

\section{Critical Issues}

\subsection{Redundant Pipeline and Descriptor Set Bindings (Critical)}

\begin{tabular}{ll}
\textbf{Severity} & \severity{Critical} \\
\textbf{Files}     & All renderers in \texttt{vulkan/renderer/} \\
\textbf{Pattern}   & Multiple \texttt{vkCmdBindPipeline} and
                      \texttt{vkCmdBindDescriptorSets} per pass \\
\end{tabular}

\vspace{0.5em}
\noindent
\textbf{Description:} Renderers issue \texttt{vkCmdBindPipeline} and
\texttt{vkCmdBindDescriptorSets} \textbf{before every draw call} rather than
checking if the pipeline or descriptor set has actually changed. For example,
\texttt{SolidRenderer} binds the pipeline, descriptor set, and push constants
for each indirect draw command in the submitted batch --- even when the same
pipeline is used for consecutive draws. With thousands of draw calls per frame,
this produces tens of thousands of redundant Vulkan API calls that the driver
must validate and process.

\noindent\textbf{Recommendation:} Track the currently bound pipeline and
descriptor set per command buffer. Skip binding calls when the state has not
changed. Group consecutive draws that share the same pipeline. This is a
mechanical change that eliminates driver overhead proportional to draw count.

\noindent\textbf{Estimated improvement:} $50{-}80\%$ reduction in CPU-side
binding API calls, saving $100{-}300\us$ per frame in command buffer recording.

\subsection{Per-Frame Descriptor Set Allocation (Critical)}

\begin{tabular}{ll}
\textbf{Severity} & \severity{Critical} \\
\textbf{Files}     & \texttt{VulkanApp.cpp}, \texttt{SceneRenderer.cpp},
                      \texttt{VegetationRenderer.cpp} \\
\textbf{Pattern}   & \texttt{vkCreateDescriptorPool}/\texttt{vkCreateDescriptorSet}
                      per frame \\
\end{tabular}

\vspace{0.5em}
\noindent
\textbf{Description:} Multiple code paths allocate fresh descriptor sets every
frame instead of recycling from a pre-allocated pool:

\begin{itemize}[nosep]
    \item \texttt{SceneRenderer} allocates descriptor sets for texture array
          bindings each frame.
    \item \texttt{VegetationRenderer} allocates per-chunk descriptor sets for
          compute dispatches.
    \item The async 360\degree{} cubemap and back-face passes create descriptor
          pools and sets per iteration.
\end{itemize}

Each \texttt{vkCreateDescriptorPool} involves kernel-mode allocation in the
Vulkan driver. With 3 frames in flight and multiple allocators, this generates
unnecessary driver object churn.

\noindent\textbf{Recommendation:} Pre-allocate descriptor pools at startup with
\texttt{VK\_DESCRIPTOR\_POOL\_CREATE\_FREE\_DESCRIPTOR\_SET\_BIT}. Create a
ring of descriptor sets (one per frame in flight) and update them with
\texttt{vkUpdateDescriptorSets} as needed. Reset pools with
\texttt{vkResetDescriptorPool} instead of destroy/create.

\noindent\textbf{Estimated improvement:} $50{-}150\us$ per frame saved in
descriptor object creation/cleanup; reduces CPU-side frame-time variance.

\subsection{Scattered Small Memory Allocations (Critical)}

\begin{tabular}{ll}
\textbf{Severity} & \severity{Critical} \\
\textbf{Files}     & \texttt{VulkanApp.cpp}, \texttt{VulkanResourceManager.cpp},
                      all renderers \\
\textbf{Pattern}   & \texttt{vkAllocateMemory} for individual images and buffers \\
\end{tabular}

\vspace{0.5em}
\noindent
\textbf{Description:} The codebase makes approximately \textbf{30+ direct}
\texttt{vkAllocateMemory} calls (see report 04). Each allocation is a separate
Vulkan device memory object, each requiring at least 4 KiB
page-aligned GPU address space. Many are for small or medium-sized resources
(textures of $1024 \times 1024 \times 4$ bytes $\approx$ 4 MiB each, staging
buffers of 4 MiB, depth targets of $\sim$8 MiB). Each allocation pays the
driver overhead of a kernel call and tracks a separate \texttt{VkDeviceMemory}
handle.

The worst offenders are:
\begin{itemize}[nosep]
    \item \textbf{Per-pass depth/color targets}: each renderer (solid, water,
          shadow $\times$ 3 cascades, back-face, 360\degree) creates its own
          depth and color attachments as separate allocations.
    \item \textbf{Staging buffers}: when the ring buffer is exhausted, new
          staging buffers are allocated on demand.
    \item \textbf{Texture arrays}: 5 arrays $\times$ 32 layers each, allocated
          individually or in small groups.
\end{itemize}

\noindent\textbf{Recommendation:} Adopt Vulkan Memory Allocator (VMA) for all
image allocations, as is already done for buffers (see
\texttt{VulkanApp.cpp:4305}). VMA will sub-allocate from large heaps,
dramatically reducing the number of \texttt{vkAllocateMemory} calls. For the
texture arrays, consider allocating a single large image with array layers
instead of separate images.

\noindent\textbf{Estimated improvement:} Reduces \texttt{vkAllocateMemory} calls
from 30+ to 2--5 per initialization. Saves $5{-}15\%$ GPU memory through
reduced alignment waste.

\section{High-Severity Issues}

\subsection{No Pipeline Caching on Disk (High)}

\begin{tabular}{ll}
\textbf{Severity} & \severity{High} \\
\textbf{Files}     & Entire codebase \\
\textbf{Pattern}   & \texttt{VK\_NULL\_HANDLE} for \texttt{VkPipelineCache} \\
\end{tabular}

\vspace{0.5em}
\noindent
\textbf{Description:} Every pipeline creation call uses
\texttt{VK\_NULL\_HANDLE} as the cache parameter. The engine creates
approximately 20--30 graphics pipelines and 5--10 compute pipelines at
startup. Each pipeline compiles SPIR-V to GPU machine code --- a process that
can take hundreds of milliseconds per pipeline on first creation. With no disk
cache, every application run repeats this compilation.

\noindent\textbf{Recommendation:} Create a \texttt{VkPipelineCache} at
initialization. Serialize to disk on shutdown via
\texttt{vkGetPipelineCacheData}. Reload at startup if the file exists. Target
\texttt{VK\_KHR\_pipeline\_binary} (Vulkan 1.4) for granular cache
invalidation when available.

\noindent\textbf{Estimated improvement:} Reduces startup time by
$1{-}3$ seconds on cold start; near-instant on subsequent runs.

\subsection{Implicit Layout Transitions via Render Pass (High)}

\begin{tabular}{ll}
\textbf{Severity} & \severity{High} \\
\textbf{Files}     & \texttt{ShadowRenderer.cpp}, \texttt{SolidRenderer.cpp},
                      \texttt{WaterRenderer.cpp} \\
\textbf{Pattern}   & \texttt{VK\_ATTACHMENT\_LOAD\_OP\_CLEAR}/\texttt{STORE\_OP\_STORE}
                      with implicit transitions \\
\end{tabular}

\vspace{0.5em}
\noindent
\textbf{Description:} While the engine uses dynamic rendering, several renderers
rely on \texttt{vkCmdBeginRendering} to handle image layout transitions
implicitly. This means the driver must insert pipeline barriers based on the
initial and final layouts specified in \texttt{VkRenderingAttachmentInfo}.
When the actual image layout at the time of the call does not match the
declared initial layout, the driver silently inserts a full
\texttt{TOP\_OF\_PIPE} to \texttt{BOTTOM\_OF\_PIPE} barrier, eliminating any
fine-grained synchronization.

Specific instances where this mismatch is likely:
\begin{itemize}[nosep]
    \item Shadow cascade depth images are rendered, then sampled in the next
          pass. The transition from \texttt{DEPTH\_ATTACHMENT\_OPTIMAL} to
          \texttt{SHADER\_READ\_ONLY\_OPTIMAL} is driver-managed.
    \item Water offscreen target transitions between color attachment and
          shader read layouts.
\end{itemize}

\noindent\textbf{Recommendation:} Manage all image layout transitions explicitly
with \texttt{vkCmdPipelineBarrier2} rather than relying on
\texttt{VkRenderingAttachmentInfo} initial/final layouts. Use the
\texttt{VK\_ATTACHMENT\_LOAD\_OP\_LOAD} pattern with explicit barriers for
subpass inputs. Document each barrier with its purpose.

\noindent\textbf{Estimated improvement:} Eliminates implicit full-pipeline
barriers, potentially saving $0.2{-}0.5\ms$ per frame in driver overhead.

\subsection{Push Constant Overuse (High)}

\begin{tabular}{ll}
\textbf{Severity} & \severity{High} \\
\textbf{Files}     & \texttt{shaders/vegetation.vert},
                      \texttt{shaders/main.frag},
                      \texttt{vulkan/renderer/SceneRenderer.cpp},
                      \texttt{vulkan/renderer/IndirectRenderer.cpp} \\
\textbf{Pattern}   & Push constants > 128 bytes; frequent \texttt{vkCmdPushConstants} \\
\end{tabular}

\vspace{0.5em}
\noindent
\textbf{Description:} Push constants are used extensively across all renderers.
While convenient, the Vulkan minimum guaranteed limit is 128 bytes, and some
shaders push more than this (vegetation wind at $\sim$160 bytes). Beyond the
portability concern, \texttt{vkCmdPushConstants} has driver cost: on some
implementations, flushing push constants can stall the command processor.

Key offenders:
\begin{itemize}[nosep]
    \item \textbf{Vegetation wind parameters}: $\sim$160 bytes (exceeds
          minimum spec).
    \item \textbf{Indirect culling}: 72-byte push constant per dispatch,
          called once per renderer per frame.
    \item \textbf{Per-draw push constants}: issued for every draw call in
          several renderers, even when the data is static across multiple draws.
\end{itemize}

\noindent\textbf{Recommendation:} Keep push constants under 128 bytes (Vulkan
guaranteed limit). Move large or infrequently changing data to UBOs. For
per-draw data, consider using a vertex buffer or storage buffer with
\texttt{gl\_VertexIndex} / \texttt{gl\_DrawID} indexing.

\noindent\textbf{Estimated improvement:} $30{-}80\us$ saved per frame in
push-constant-related driver overhead; ensures portability to mobile/integrated
GPUs.

\subsection{Command Buffer Reset vs.\ Reallocate (High)}

\begin{tabular}{ll}
\textbf{Severity} & \severity{High} \\
\textbf{Files}     & \texttt{VulkanApp.cpp}, \texttt{VegetationRenderer.cpp},
                      \texttt{IndirectRenderer.cpp} \\
\textbf{Pattern}   & \texttt{vkFreeCommandBuffers} + \texttt{vkAllocateCommandBuffers}
                      per frame \\
\end{tabular}

\vspace{0.5em}
\noindent
\textbf{Description:} Several code paths destroy and reallocate command buffers
every frame instead of resetting them with \texttt{vkResetCommandBuffer}. The
allocate/free cycle is more expensive than reset because it involves the
driver's allocation tracking. The worst instance is the async 360\degree{}
cubemap and back-face passes, which allocate and free their command buffers
every iteration.

\noindent\textbf{Recommendation:} Allocate command buffers once at
initialization (or pool creation) and use \texttt{vkResetCommandBuffer} with
\texttt{VK\_COMMAND\_BUFFER\_RESET\_RELEASE\_RESOURCES\_BIT} as needed. Prefer
\texttt{vkResetCommandPool} for bulk resets.

\noindent\textbf{Estimated improvement:} $20{-}50\us$ per frame saved in
command buffer allocation overhead.

\subsection{Per-Frame \texttt{vkCmdSetDepthBounds} Calls (High)}

\begin{tabular}{ll}
\textbf{Severity} & \severity{High} \\
\textbf{Files}     & \texttt{ShadowRenderer.cpp},
                      \texttt{SolidRenderer.cpp} \\
\textbf{Pattern}   & \texttt{vkCmdSetDepthBounds} called every frame with
                      (0.0, 1.0) \\
\end{tabular}

\vspace{0.5em}
\noindent
\textbf{Description:} Several renderers call \texttt{vkCmdSetDepthBounds(0.0f,
1.0f)} every frame, even when depth bounds testing is not enabled in the
pipeline. This is a redundant state command that the driver must process.
Additionally, the pipeline is created with
\texttt{VK\_DYNAMIC\_STATE\_DEPTH\_BOUNDS}, which requires this call even
when the feature is unused.

\noindent\textbf{Recommendation:} Remove \texttt{VK\_DYNAMIC\_STATE\_DEPTH\_BOUNDS}
from pipelines that do not use depth bounds testing. Remove the corresponding
\texttt{vkCmdSetDepthBounds} call. This eliminates a redundant driver-level
state update.

\noindent\textbf{Estimated improvement:} $5{-}15\us$ per frame saved in driver
state processing.

\subsection{Shader Module Creation/Destruction Per Run (High)}

\begin{tabular}{ll}
\textbf{Severity} & \severity{High} \\
\textbf{Files}     & \texttt{VulkanApp.cpp}, \texttt{VulkanResourceManager.cpp} \\
\textbf{Pattern}   & \texttt{vkCreateShaderModule} for every pipeline at
                      startup \\
\end{tabular}

\vspace{0.5em}
\noindent
\textbf{Description:} Shader modules are created from SPIR-V at every
application launch and destroyed at shutdown. While this is expected, the
engine does not retain shader modules for reuse across pipelines that share
the same shader stage. Each pipeline creation for a different specialization
constant combination re-reads the SPIR-V binary from disk and creates a new
module. With 20--30 pipelines, this multiplies the I/O and module creation
cost.

\noindent\textbf{Recommendation:} Load SPIR-V once per shader file and keep the
\texttt{VkShaderModule} alive for the application lifetime. Reuse across all
pipelines. With Vulkan 1.3+, consider using \texttt{VK\_EXT\_shader\_module\_identifier}
for even faster pipeline creation.

\noindent\textbf{Estimated improvement:} Reduces startup time by
$200{-}500\ms$ by eliminating redundant SPIR-V loads and module creations.

\section{Medium-Severity Issues}

\subsection{No Clear Color/Depth Optimisation (Medium)}

\begin{tabular}{ll}
\textbf{Severity} & \severity{Medium} \\
\textbf{Files}     & All renderers with \texttt{LOAD\_OP\_CLEAR} \\
\textbf{Pattern}   & \texttt{VK\_ATTACHMENT\_LOAD\_OP\_CLEAR} on all targets \\
\end{tabular}

\vspace{0.5em}
\noindent
\textbf{Description:} Every offscreen target uses
\texttt{VK\_ATTACHMENT\_LOAD\_OP\_CLEAR} for color and depth attachments,
even when the entire attachment will be overwritten (e.g.\ solid color pass
after depth pre-pass, where every pixel is shaded). This forces a clear
operation that may be redundant when the subsequent draw covers the full
attachment.

\noindent\textbf{Recommendation:} Use \texttt{VK\_ATTACHMENT\_LOAD\_OP\_DONT\_CARE}
for attachments that are fully covered by subsequent rendering. This allows
the driver to skip the clear. For the depth pre-pass, \texttt{DONT\_CARE}
is always safe because every pixel written contains a valid depth value.

\noindent\textbf{Estimated improvement:} $10{-}30\us$ per cleared attachment
per frame. With 5--8 attachments per frame, saves $50{-}200\us$ total.

\subsection{Single-Threaded Resource Loading (Medium)}

\begin{tabular}{ll}
\textbf{Severity} & \severity{Medium} \\
\textbf{Files}     & \texttt{main.cpp:setupTextures()},
                      \texttt{main.cpp:186--383} \\
\textbf{Pattern}   & Sequential file I/O + decode on main thread \\
\end{tabular}

\vspace{0.5em}
\noindent
\textbf{Description:} All 100+ texture images (20 texture triples) are loaded
sequentially on the main thread during initialization. Each file involves:
disk read (several syscalls), JPEG decompression (stb\_image), sRGB to linear
conversion, and GPU upload. With no parallel loading, this adds $1{-}3$
seconds to startup on HDDs and $\sim$0.5 seconds on SSDs.

\noindent\textbf{Recommendation:} Use \texttt{std::async} or the existing
thread pool to decode textures in parallel. Upload to GPU via staging buffer
on background threads. Show a loading screen with progress indication.

\noindent\textbf{Estimated improvement:} Reduces startup time by
$30{-}70\%$ depending on storage speed and CPU core count.

\subsection{No Dedicated Transfer Queue Usage (Medium)}

\begin{tabular}{ll}
\textbf{Severity} & \severity{Medium} \\
\textbf{Files}     & \texttt{VulkanApp.cpp} \\
\textbf{Pattern}   & \texttt{VK\_QUEUE\_TRANSFER\_BIT} queue exists but is
                      rarely used \\
\end{tabular}

\vspace{0.5em}
\noindent
\textbf{Description:} The engine queries for a dedicated transfer queue during
device enumeration but uses it only sporadically. Most buffer and image uploads
go through the graphics queue, competing with rendering work. A dedicated
transfer queue allows the driver to overlap data movement with rendering.

\noindent\textbf{Recommendation:} Route all staging buffer copies, texture
uploads, and mesh data transfers through the dedicated transfer queue when
available. Use timeline semaphores to signal completion to the graphics queue.

\noindent\textbf{Estimated improvement:} $100{-}300\us$ per frame saved during
upload-heavy frames by overlapping DMA with rendering.

\subsection{Per-Frame Indirect Draw Count Readback (Medium)}

\begin{tabular}{ll}
\textbf{Severity} & \severity{Medium} \\
\textbf{Files}     & \texttt{vulkan/renderer/IndirectRenderer.cpp},
                      \texttt{main.cpp} \\
\textbf{Pattern}   & Host read of \texttt{VK\_BUFFER\_USAGE\_INDIRECT\_BUFFER\_BIT}
                      content \\
\end{tabular}

\vspace{0.5em}
\noindent
\textbf{Description:} The ImGui stats overlay reads visible draw counts by
mapping and reading the indirect buffer that was written by the GPU culling
compute shader. This read requires a CPU--GPU synchronization (either an
implicit barrier or a fence wait) and forces the GPU to flush its caches.
Even when the stats overlay is hidden, the read may still occur.

\noindent\textbf{Recommendation:} Cache visible counts in a separate
\texttt{HOST\_COHERENT} buffer that the compute shader also writes to. Read
it only when the stats overlay is visible. Batch reads to once every 10--30
frames.

\noindent\textbf{Estimated improvement:} $30{-}60\us$ per frame when stats are
hidden; eliminates a potential GPU stall.

\subsection{Per-Frame Texture Array Descriptor Updates (Medium)}

\begin{tabular}{ll}
\textbf{Severity} & \severity{Medium} \\
\textbf{Files}     & \texttt{SceneRenderer.cpp:updateTextureDescriptorSet()} \\
\textbf{Pattern}   & Full descriptor set update every frame \\
\end{tabular}

\vspace{0.5em}
\noindent
\textbf{Description:} The texture array descriptor set is fully updated every
frame with \texttt{vkUpdateDescriptorSets}, even though the actual texture
contents rarely change (only during TextureMixer generation). Most frames
simply rebind the same 5 arrays (albedo, normal, bump, roughness, AO) via
descriptor writes.

\noindent\textbf{Recommendation:} Update the descriptor set only when the
texture layout or binding actually changes. Use a dirty flag to track whether
the texture arrays have been modified since the last update.

\noindent\textbf{Estimated improvement:} $10{-}30\us$ per frame saved in
descriptor update overhead.

\subsection{ImGui Continuous Stats Polling (Medium)}

\begin{tabular}{ll}
\textbf{Severity} & \severity{Medium} \\
\textbf{Files}     & \texttt{main.cpp:drawSidebar()} \\
\textbf{Pattern}   & Per-frame queries to \texttt{getMeshCount()},
                      \texttt{getVegetationCount()}, etc. \\
\end{tabular}

\vspace{0.5em}
\noindent
\textbf{Description:} The ImGui stats window calls multiple accessor functions
every frame that acquire locks, iterate over containers, or read GPU buffers.
These include the visible count reads mentioned above, plus mesh counts from
\texttt{IndirectRenderer} (under mutex) and vegetation chunk counts. When the
stats overlay is not visible, these calls still execute.

\noindent\textbf{Recommendation:} Guard all stats queries with an
\texttt{ImGui::IsItemVisible()} or a global \texttt{showStats} flag. Cache
aggregate values for 10--30 frames when the overlay is visible.

\noindent\textbf{Estimated improvement:} $30{-}100\us$ per frame saved when
stats overlay is hidden.

\subsection{No Relaxed Line Ordering in Descriptor Writes (Medium)}

\begin{tabular}{ll}
\textbf{Severity} & \severity{Medium} \\
\textbf{Files}     & \texttt{SceneRenderer.cpp}, \texttt{VulkanApp.cpp} \\
\textbf{Pattern}   & Sequential \texttt{VkWriteDescriptorSet} entries \\
\end{tabular}

\vspace{0.5em}
\noindent
\textbf{Description:} When updating descriptor sets with multiple image or
buffer bindings, the code creates individual \texttt{VkWriteDescriptorSet}
on the stack and calls \texttt{vkUpdateDescriptorSets} once with an array.
This is correct but creates per-frame stack allocation and setup overhead
for the array and each entry.

\noindent\textbf{Recommendation:} Store the \texttt{VkWriteDescriptorSet}
templates as static const arrays or member variables. Only update the
changing fields (e.g.\ buffer handles, image views) each frame using
assignment rather than reconstruction.

\noindent\textbf{Estimated improvement:} $5{-}15\us$ per frame saved in
descriptor write setup.

\section{Low-Severity Issues}

\subsection{Unused Shader Module Retention (Low)}

\begin{tabular}{ll}
\textbf{Severity} & \severity{Low} \\
\textbf{Files}     & \texttt{VulkanApp.cpp} \\
\textbf{Pattern}   & Shader modules kept alive after pipeline creation \\
\end{tabular}

\vspace{0.5em}
\noindent
\textbf{Description:} Shader modules are retained for the entire application
lifetime even after all pipelines referencing them have been created. While
this is harmless (modules are small), destroying them after pipeline creation
would free minor driver resources.

\noindent\textbf{Recommendation:} Destroy shader modules immediately after
pipeline creation when using \texttt{VK\_EXT\_shader\_module\_identifier}
(Vulkan 1.3) or when the module is no longer needed for further pipeline
creations.

\subsection{Duplicate \texttt{vkCmdSetViewport} Calls (Low)}

\begin{tabular}{ll}
\textbf{Severity} & \severity{Low} \\
\textbf{Files}     & \texttt{SolidRenderer.cpp}, \texttt{WaterRenderer.cpp},
                      \texttt{ShadowRenderer.cpp} \\
\textbf{Pattern}   & \texttt{vkCmdSetViewport} called in every render pass \\
\end{tabular}

\vspace{0.5em}
\noindent
\textbf{Description:} Each renderer calls \texttt{vkCmdSetViewport} with the
same full-screen viewport at the start of its command buffer. Since the
viewport rarely changes between passes, these calls are redundant. This is
a very minor overhead.

\noindent\textbf{Recommendation:} Set the viewport once at the start of the
main command buffer and only reissue when the render area changes (e.g.\
shadow cascade resolution changes).

\subsection{Per-Frame \texttt{vkCmdSetScissor} Redundancy (Low)}

\begin{tabular}{ll}
\textbf{Severity} & \severity{Low} \\
\textbf{Files}     & \texttt{SolidRenderer.cpp}, \texttt{WaterRenderer.cpp} \\
\textbf{Pattern}   & \texttt{vkCmdSetScissor} set per render pass \\
\end{tabular}

\vspace{0.5em}
\noindent
\textbf{Description:} Same pattern as viewport: scissor rect is set to the
full attachment size in every pass, even when it does not change between
passes. The driver must process each \texttt{vkCmdSetScissor} command.

\noindent\textbf{Recommendation:} Set scissor once per frame, similar to
viewport recommendation.

\subsection{No Mipmap Generation During Load (Low)}

\begin{tabular}{ll}
\textbf{Severity} & \severity{Low} \\
\textbf{Files}     & \texttt{TextureArrayManager.cpp} \\
\textbf{Pattern}   & \texttt{mipLevels = 1} for all texture arrays \\
\end{tabular}

\vspace{0.5em}
\noindent
\textbf{Description:} Noted in previous reports (report 02) but still present:
texture arrays are created with a single mip level. While fixing this requires
texture generation infrastructure, the lack of mipmaps reduces texture cache
efficiency when sampling at a distance.

\noindent\textbf{Recommendation:} Generate mip chains at texture array creation
time using \texttt{vkCmdBlitImage}. Update the fragment shader to use
\texttt{textureLod} with computed LOD or rely on hardware auto-mip selection.

\section{Pipeline State Management}

\subsection{Pipeline Count and Specialization Constants}

The engine creates graphics pipelines for each renderer with unique
specialization constant combinations:

\begin{center}
\begin{tabular}{lccc}
\toprule
\textbf{Renderer} & \textbf{Pipelines} & \textbf{Spec Constants} & \textbf{Notes} \\
\midrule
SolidRenderer & 2 & shadow cascade level & Depth + color variants \\
ShadowRenderer & 1 & -- & Single depth pipeline \\
SkyRenderer & 1 & -- & Full-screen triangle \\
VegetationRenderer & 3 & LOD, billboard mode & Depth, color, impostor \\
WaterRenderer & 2 & refraction flag & Geometry + back-face \\
IndirectRenderer & 2 & -- & Cull, sort computes \\
ImpostorCapture & 2 & face index & 6 faces of cubemap \\
DebugCubeRenderer & 1 & -- & Debug overlay \\
PostProcessRenderer & 1 & -- & Final composite \\
\bottomrule
\end{tabular}
\end{center}

Each pipeline is created with \texttt{VK\_NULL\_HANDLE} cache, meaning all
15+ pipelines recompile every run.

\subsection{Descriptor Set Layout Count}

The engine uses the following distinct descriptor set layouts:

\begin{center}
\begin{tabular}{lcc}
\toprule
\textbf{Layout} & \textbf{Bindings} & \textbf{Used By} \\
\midrule
Main UBO & 0: UBO\$1 & All renderers \\
Texture arrays & 0--4: combined image samplers & Solid, vegetation \\
Shadow samplers & 0: sampler, 1: sampled image & Solid, water \\
Water UBO & 0: UBO & WaterRenderer \\
Sky UBO & 0: UBO, 1: sampled image & SkyRenderer \\
Cube environment & 0: combined image sampler & Solid (reflections) \\
Vegetation instance & 0: SSBO & VegetationRenderer compute \\
Indirect cull & 0--3: SSBO \& UBO & IndirectRenderer compute \\
\bottomrule
\end{tabular}
\end{center}

Total: 8 distinct layouts. This is manageable but means descriptor set
switching between renderers is frequent.

\section{Resource Lifecycle Analysis}

\subsection{Creation vs.\ Destruction Order}

The engine follows RAII principles but does not rigorously enforce reverse
creation-order destruction in all paths. Key observations:

\begin{itemize}[nosep]
    \item \textbf{Swapchain}: Recreated correctly with proper drain before
          destruction.
    \item \textbf{Descriptor pools}: Many are allocated per-frame and freed
          via \texttt{deferDestroyUntilFence}. This works but creates churn.
    \item \textbf{Command pools}: Similar per-frame allocation pattern.
    \item \textbf{Buffers}: Persistent buffers (UBO, SSBO) are created once
          and destroyed at shutdown. Transient staging buffers are handled
          via \texttt{deferDestroyUntilFence}.
    \item \textbf{Images}: Most images persist for the application lifetime
          (textures, depth targets). Some temporary images (e.g.\ 360\degree
          cubemap captures) are created and destroyed per-frame.
\end{itemize}

\subsection{Deferred Destruction Queue Pressure}

The \texttt{deferDestroyUntilFence} mechanism is used extensively. At any
point, dozens of Vulkan objects may be pending destruction. The pending
queue is checked every frame via \texttt{vkGetFenceStatus}, which polls the
GPU state. While this avoids blocking, it does add per-frame overhead.

The queue typically contains:
\begin{itemize}[nosep]
    \item 2--3 descriptor pools (from async passes)
    \item 4--6 descriptor sets
    \item 2--3 command buffers
    \item 1--2 staging buffers
    \item 1--2 fences
    \item 0--2 semaphores
\end{itemize}

This is 10--18 objects pending destruction at any time. The overhead of
tracking and processing these is small but non-zero.

\section{Summary and Global Scoring}

\subsection{Issue Summary}

\begin{longtable}{p{0.30\textwidth}cp{0.20\textwidth}p{0.20\textwidth}}
    \toprule
    \textbf{Issue} & \textbf{Sev.} & \textbf{Effort} & \textbf{Impact} \\
    \midrule
    \endhead
    \bottomrule
    \endfoot
    Redundant pipeline/desc.\ binds
        & \severity{Critical}
        & Medium
        & $100{-}300\us$ CPU saved \\
    Per-frame descriptor allocation
        & \severity{Critical}
        & Medium
        & $50{-}150\us$ CPU saved \\
    Scattered small allocations
        & \severity{Critical}
        & High
        & $5{-}15\%$ memory saved \\
    No pipeline disk cache
        & \severity{High}
        & Low
        & $1{-}3$s startup saved \\
    Implicit layout transitions
        & \severity{High}
        & Medium
        & $0.2{-}0.5\ms$ GPU saved \\
    Push constant overuse
        & \severity{High}
        & Medium
        & Portability + $30{-}80\us$ \\
    Cmd buffer reset vs.\ free
        & \severity{High}
        & Low
        & $20{-}50\us$ CPU saved \\
    Depth bounds redundancy
        & \severity{High}
        & Trivial
        & $5{-}15\us$ CPU saved \\
    Shader module churn
        & \severity{High}
        & Low
        & $200{-}500\ms$ startup \\
    Clear colour optimisation
        & \severity{Medium}
        & Low
        & $50{-}200\us$ GPU saved \\
    Single-threaded loading
        & \severity{Medium}
        & Medium
        & $30{-}70\%$ startup \\
    Transfer queue underuse
        & \severity{Medium}
        & Medium
        & $100{-}300\us$ upload \\
    Draw count readback
        & \severity{Medium}
        & Low
        & $30{-}60\us$ stall \\
    Texture descriptor updates
        & \severity{Medium}
        & Low
        & $10{-}30\us$ saved \\
    ImGui stats polling
        & \severity{Medium}
        & Low
        & $30{-}100\us$ saved \\
    Descriptor write setup
        & \severity{Medium}
        & Trivial
        & $5{-}15\us$ saved \\
    Shader module retention
        & \severity{Low}
        & Trivial
        & Minor \\
    Viewport/scissor redundancy
        & \severity{Low}
        & Trivial
        & $5{-}10\us$ saved \\
    No mipmaps
        & \severity{Low}
        & Medium
        & Texture cache \\
\hline
\end{longtable}

\subsection{Global Score}

The engine shows solid architectural choices (dynamic rendering, GPU culling,
async compute) but loses performance in the details of CPU-side recording
efficiency and resource management.

\begin{center}
\begin{tabular}{lcc}
\toprule
\textbf{Category} & \textbf{Score} & \textbf{Max} \\
\midrule
Architecture & 8/10 & Modern, dynamic rendering, GPU-driven \\
CPU efficiency & 5/10 & Redundant binds, per-frame allocations \\
GPU efficiency & 6/10 & No mipmaps, implicit barriers \\
Memory management & 4/10 & Scattered allocations, no VMA images \\
Synchronization & 6/10 & Correct but heavy-handed \\
Resource lifecycle & 5/10 & Per-frame create/destroy churn \\
Build \& startup & 4/10 & No pipeline cache, serial loading \\
Portability & 6/10 & Push constants > 128B \\
\midrule
\textbf{Overall} & \textbf{5.5/10} & \textbf{Fair} \\
\bottomrule
\end{tabular}
\end{center}

\section{Immediate Recommendations}

Based on effort/impact ratio:

\begin{enumerate}
    \item \textbf{Track bound pipeline and descriptor set state per command
          buffer} (Section 5.1) --- medium effort, eliminates majority of
          redundant driver calls.
    \item \textbf{Pre-allocate descriptor pools and sets} (Section 5.2) ---
          medium effort, eliminates per-frame Vulkan object churn.
    \item \textbf{Add a pipeline cache with disk serialization}
          (Section 6.1) --- low effort, immediate startup time improvement.
    \item \textbf{Replace command buffer free/allocate with reset}
          (Section 6.4) --- low effort, reduces driver overhead.
    \item \textbf{Remove redundant \texttt{vkCmdSetDepthBounds} and dynamic
          state from pipelines that do not use it} (Section 6.5) --- trivial
          effort.
    \item \textbf{Adopt VMA for image allocations} (Section 5.3) --- high
          effort but transformative for memory efficiency.
    \item \textbf{Load shaders once and reuse modules} (Section 6.6) ---
          low effort, measurable startup improvement.
\end{enumerate}

\section{Conclusion}

Performance report \#05 identifies a central theme: the engine is functionally
correct and architecturally modern, but pays a significant CPU tax through
redundant API calls and per-frame object creation. The single biggest gain
comes from tracking GPU state during command buffer recording to eliminate
unnecessary \texttt{vkCmdBind*} calls. Combined with pipeline caching and
descriptor set pre-allocation, these changes could reduce per-frame CPU
overhead by $200{-}500\us$ --- a meaningful improvement for scenes that are
currently CPU-bound.

The overall rating remains \textbf{Fair (5.5/10)}, reflecting strengths in
architecture and GPU-side design alongside weaknesses in CPU-side recording
efficiency and resource management. Addressing the Critical and High issues
identified here would bring the engine to a \textbf{Good (7/10)} rating.

\end{document}
